% W745: proposed replacement abstract for tnf_paper.tex, under the Nine-Rung Law.
% NOT applied to the paper. Review, then splice between \begin{abstract} and
% \end{abstract}. Every number below has a run behind it; provenance is listed
% at the end of this file.

\begin{abstract}
A ternary node has three sites that could carry a number format; one needs to.
The weight is a code: multiplying by it is a sign-select. The sample arrives
ADC-native. The accumulator alone has a range to spend, and Ternary Network
Floats spend it --- a sign, a balanced-ternary exponent of $E_t$ trits, an
$M$-bit mantissa, $1+E_t+M=N$.

\paragraph{Budget convention, stated once and binding throughout.}\label{sec:budget}
$1+E_t+M=N$ counts the exponent as $E_t$ \emph{trit positions}. It is a
specification convention, not a storage claim, and this paper uses two budgets
under two explicit names. The \textbf{nominal} budget is $1+E_t+M=N$; the
\textbf{physical} budget is $1+\lceil E_t\log_2 3\rceil+M \le N$, which is what a
packed binary word costs. The two differ by
$\lceil E_t\log_2 3\rceil - E_t$ cells, three at $E_t=4$ and seven at $E_t=11$, so
the adopted 16-bit rung $(E_t,M)=(4,11)$ occupies $19$ physical bits. Every
accuracy comparison against a binary format is stated with its budget named. Under
the physical budget TNF holds no advantage at $16$ or $32$ bits ---
Section~\ref{sec:threeresults} measures this, Table~\ref{tab:accphys} states it
for the accuracy table specifically, and Theorem~\ref{thm:rangeperbit} says why it
cannot be otherwise. Claims of the form ``TNF wins at equal width'' appear nowhere
in this paper without the budget attached.

\paragraph{The Nine-Rung Law, and what it replaces.}\label{sec:ninerung}
The weight half of this paper began as a claim about a \emph{set}: that the
alphabet $\{-\varphi,0,+\varphi\}$ is where a multiplier is decided, because
$\mathbb{Z}[\varphi]$ is closed under weight application and accumulation alike
and $\varphi^2=\varphi+1$ makes a weight the Fibonacci step
$(a,b)\mapsto(b,a+b)$. That closure is proved and machine-checked in Coq, and it
is retained below. \textbf{The claim built on it is withdrawn.} To test it we
constructed the \emph{golden ladder}
$\mathrm{GA\text{-}T}_n=\{0\}\cup\{\pm\varphi^k: 0\le k\le n\}$, of cardinality
$2n{+}3$, and measured it on five axes. The ladder answered against itself, and
the measurement it made possible is the result this paper now reports
\textup{[measured]}:

\begin{enumerate}
\item \textbf{Cardinality dominates shape by an order of magnitude.} Across three
binarised tasks (UNSW-NB15, MNIST, Fashion-MNIST), seven alphabets and $30$ seeds
each --- $630$ trained runs, paired $t$-tests, inverse-variance pooled --- the
\emph{size} of the alphabet is worth $+0.844$\,pp from three levels to nine
($z=33.9$, direction identical on all three tasks), while its \emph{shape} at
fixed size is worth $+0.085$\,pp ($z=5.7$, homogeneous at seven levels,
$Q=2.2$; heterogeneous at nine, $Q=6.9$).

\item \textbf{The return saturates, and never later than the ninth rung.} Eight
rungs, three tasks, $30$ seeds --- a further $720$ runs --- then five
same-dataset tasks of graded difficulty at five rungs and $20$ seeds, a further
$500$. \textbf{Across eight tasks and $1220$ runs no step above nine levels is
significant anywhere.} The last significant rung is nine on UNSW-NB15 and
Fashion-MNIST and \emph{five} on MNIST: the ceiling is a constant, the rung is a
property of the task, and we cannot predict which: task difficulty was proposed
as the mechanism and refuted ($r=-0.35$, $n=5$, non-monotone), so the ceiling is
reported as an empirical regularity with no mechanism claimed. The first step,
three levels to five, is the largest everywhere ($+0.547$, $+1.622$, $+0.511$\,pp).

\item \textbf{Therefore the alphabet is chosen on cost, not on accuracy.} At
equal cardinality the difference between alphabets is under a tenth of a point,
and the difference in area is not: on an XC7A200T, nine dyadic levels place in
$1200$ LUT against the golden ladder's $1677$, with no DSP in either
\textup{[measured, post-route]}. A power-of-two weight is a shift into one
accumulator; a $\varphi^k$ weight is $F(k{-}1)x$ into one lane and $F(k)x$ into a
second, and pays for both. \emph{Multiplier-free is not the same as free.}
\end{enumerate}

\paragraph{GA-T is reported here as an instrument.}
The golden ladder is not this paper's answer and is the reason the paper has one.
Nothing in the dyadic result was reachable without a graded family to compare
against: the ladder is what turned ``is $\varphi$ special?'' --- unanswerable ---
into ``how much does the shape of an alphabet buy at fixed size?'', which
$1350$ runs answer. It is named, its rungs are defined by the powers of the
fundamental unit of $\mathbb{Z}[\varphi]$, its cost is measured at every rung,
and it has been placed, routed and read back off three XC7A200T dice. It is
reported in full, including the five measurements that exclude it.

\paragraph{The remaining ternary claims.}
Closure does not single out $\varphi$; enumeration does. A scale is
multiplier-free exactly when it is an algebraic integer whose companion matrix
has entries in $\{0,\pm1\}$; degree two admits $\varphi$ alone, and nothing lies
between the shift and $\varphi$. The family $r^d=r+1$ reaches any granularity at
\emph{one addition} and $d$ registers: fineness costs registers, not adders. This
is a statement about \emph{which scales are multiplier-free}, and the Nine-Rung
Law is a statement about \emph{how much any of them buys}. The first survives the
second.

On an XC7A200T the folded ternary neuron costs $28.0$ LUT per weight at fan-in
$8$, rising to $33.1$ and $33.2$ at fan-in $16$ and $32$, with \emph{no DSP at
any fan-in} (Table~\ref{tab:node}).

Applicability of the accumulator format is stated as a threshold and then tested
against it. The scalar is $D=\operatorname{mean}(\lvert\log_2|x|\rvert)$; the
threshold is a property of the rung, not of the family --- $D \gtrsim 9.5$ at
sixteen physical cells and $D \gtrsim 21.9$ at thirty-two, with a third rung
admitting no threshold at all --- and it is a conjunction with the range
condition, which does most of the excluding. Eleven intermediate quantities of a
real GPT-2 block satisfy none of it, because normalisation is what destroys the
property. Seven of twenty-two candidate pairs outside neural inference do satisfy
it, at $1.23\times$ to $4.00\times$. \texttt{takum} reaches $\pm 599$ binades at
both widths compared and never leaves range where TNF does, so the format that
wins inside the window is the one that has to be selected into it. One task
outcome rather than a round-trip error is also reported: on a
maximum-a-posteriori estimate in raw SI units with a fixed \texttt{float32}
accumulator, TNF32 lands $1.60\times$ closer to the \texttt{float64} solution than
\texttt{takum32} and settles in $21$ iterations against $27$, while
\texttt{binary32} beats both by a factor of thirty-three.

\paragraph{Retractions.}
Twenty-three retractions of our own claims are marked in place below and
enumerated in Section~\ref{sec:selfcheck}. Three are new to this revision and all
three concern the weight alphabet: that $\{-\varphi,0,+\varphi\}$ decides the
multiplier question, that the golden ladder's representation efficiency predicts
its trained accuracy, and that a dyadic advantage measured on one task
generalises. Each was refuted by the experiment proposed to confirm it.
\end{abstract}

% ---------------------------------------------------------------------------
% PROVENANCE, so every number above can be checked rather than trusted.
%
%  +0.844 pp / z=33.9   size effect, pooled over 3 tasks        T286
%  +0.085 pp / z=5.7    shape effect, pooled; Q=2.2 at K=7      T286
%  9 / 5 / 9            last significant rung per task          T288
%  +0.547/+1.622/+0.511 the 3->5 step per task                  T288
%  1200 vs 1677 LUT     placed, XC7A200T, one 64->8 layer       T280, T283
%  630 + 720 runs       3 tasks x 7 arms x 30, then 8 rungs     T286, T288
%  silicon readback     Done 0x0 -> done 1, magic 0xA5A5A5A     T219, T281
%
% All in docs/theory/IGLA-FORMAL-RESULTS.md of gHashTag/t27, branch
% claude/igla-fpga-improvements-3f5e1a. Raw runs in
% experiments/gfternary-line/*.json.
%
% NOT APPLIED and NOT PUSHED: the handoff rule is that Dmitrii confirms before
% any push or merge. The title change to TNF-9 is proposed separately and is not
% made here, because the \title block carries the canon plate and the logo.
% ---------------------------------------------------------------------------
